% =============================================================================
% APPENDIX L: STRONG CP — THEOREM-GRADE ALL-ORDERS CLOSURE
% =============================================================================
% This appendix proves that θ̄ = 0 to all loop orders in the DFD microsector.
% The key insight: the mapping torus has even dimension (8), forcing η = 0.
% =============================================================================

\section{Strong CP: All-Orders Closure via CP Non-Anomaly}\label{app:strongcp_selection_rule}

\subsection{What must be shown}

In any 4D gauge theory with quarks, the physical strong-CP parameter is
\begin{equation}
\bar\theta\;=\;\theta_{\rm bare}\;+\;\arg\det M_u\;+\;\arg\det M_d\,.
\label{eq:L_theta_bar}
\end{equation}
The statement ``$\bar\theta=0$ to all orders'' is equivalent to the statement that the \emph{full quantum effective action} respects an exact CP symmetry.
Since the operator
\begin{equation}
\mathcal{O}_\theta\;\equiv\;\frac{1}{32\pi^2}\,\mathrm{Tr}(F\wedge F)
\end{equation}
changes sign under CP, any CP-invariant quantum effective action forbids a generated coefficient for $\mathcal{O}_\theta$.
Thus the all-loops claim reduces to two conditions:
\begin{enumerate}
\item \textbf{Classical CP invariance:} the microscopic action is CP invariant at $\theta_{\rm bare}=0$.
\item \textbf{No CP anomaly:} the fermion measure (determinant/Pfaffian) is invariant under CP.
\end{enumerate}
If both hold, then $\theta_{\rm bare}=0$ is protected as a selection rule and no effective $\theta$ term can be generated.

\subsection{Tree-level CP invariance (established)}

The DFD microsector on $M = \mathbb{C}P^2\times S^3$ with gauge bundle $E = \mathcal{O}(9)\oplus\mathcal{O}^{\oplus 5}$ produces:
\begin{itemize}
\item The Standard Model gauge group $G_{\rm SM} = \mathrm{SU}(3)_C\times\mathrm{SU}(2)_L\times\mathrm{U}(1)_Y$,
\item Real Yukawa eigenvalues from the K\"ahler structure,
\item $\arg\det(M_u M_d) < 10^{-19}\,\mathrm{rad}$ (verified numerically in Appendix~\ref{app:yukawa-hierarchy}),
\item Nonzero CKM CP violation ($J \neq 0$) from geometric phases.
\end{itemize}
This satisfies Condition~(1). The all-loops upgrade requires establishing Condition~(2): CP non-anomaly.

\subsection{The Dai--Freed anomaly formula}

For a discrete symmetry $\sigma$ (here $\sigma=\mathrm{CP}$), the anomaly is a $\mathrm{U}(1)$ phase given by the holonomy of the Pfaffian/determinant line bundle over background fields.
The Dai--Freed theorem~\cite{DaiFreed1994,FreedMoore2013} expresses this holonomy as an exponentiated $\eta$-invariant on the \emph{mapping torus}.

Let $M = \mathbb{C}P^2\times S^3$ be the microsector manifold with the specified Spin$^c$ structure and gauge bundle.
Define the mapping torus:
\begin{equation}
T_{\mathrm{CP}}\;\equiv\;(M\times[0,1])\,/\,(x,0)\sim(\mathrm{CP}(x),1)\,.
\label{eq:L_mapping_torus}
\end{equation}
The CP anomaly phase is then:
\begin{equation}
A_{\mathrm{CP}}\;=\;\exp\!\left(\frac{i\pi}{2}\,\eta\!\left(D_{T_{\mathrm{CP}}}\right)\right),
\label{eq:L_anomaly_phase}
\end{equation}
where $D_{T_{\mathrm{CP}}}$ is the Spin$^c$ Dirac operator on $T_{\mathrm{CP}}$ twisted by the gauge bundle, and $\eta(\cdot)$ is the APS $\eta$-invariant~\cite{DaiFreed1994}.

\textbf{Criterion.} CP is non-anomalous iff $A_{\mathrm{CP}}=1$, i.e.\ iff $\eta(D_{T_{\mathrm{CP}}})\in 4\mathbb{Z}$.

\subsection{Theorem: $\eta$ vanishes automatically in even dimensions}

\begin{theorem}[Automatic vanishing of $\eta$ in even dimensions]\label{thm:eta_even_vanishes}
Let $X$ be a \emph{closed even-dimensional} Spin$^c$ Riemannian manifold, and let
$D_E$ denote the Spin$^c$ Dirac operator on $X$ twisted by a Hermitian vector bundle $E$ with unitary connection.
Then the spectrum of $D_E$ is symmetric about $0$, hence
\begin{equation}
\eta(D_E)=0,
\end{equation}
and therefore $\exp\!\bigl(\tfrac{i\pi}{2}\eta(D_E)\bigr)=1$.
\end{theorem}

\begin{proof}
Because $\dim X$ is even, the complex spinor bundle carries a $\mathbb{Z}_2$ grading
$S=S^+\oplus S^-$ with chirality operator $\Gamma=\mathrm{diag}(+1,-1)$.
The twisted Dirac operator is \emph{odd} with respect to this grading:
\begin{equation}
\Gamma D_E \Gamma^{-1} \;=\; -\,D_E.
\label{eq:L_chirality_odd}
\end{equation}
Consequently, if $D_E\psi=\lambda\psi$ with $\lambda\neq 0$, then
$D_E(\Gamma\psi)= -\lambda(\Gamma\psi)$, and the multiplicities of $\pm\lambda$ match exactly.
Thus the $\eta$-function, defined initially for $\Re(s)\gg 0$ by
\begin{equation}
\eta(D_E,s)=\sum_{\lambda\neq 0}\mathrm{sign}(\lambda)\,|\lambda|^{-s},
\end{equation}
vanishes identically term-by-term (each $+\lambda$ cancels a $-\lambda$), and by analytic continuation $\eta(D_E)=\eta(D_E,0)=0$.
\end{proof}

\begin{corollary}[DFD Strong-CP closure]\label{cor:DFD_ACP_trivial}
The mapping torus $T_{\mathrm{CP}}$ has dimension
\begin{equation}
\dim T_{\mathrm{CP}} = \dim M + 1 = 7 + 1 = 8 \quad\text{(even)}.
\end{equation}
The CP involution on $\mathbb{C}P^2$ (complex conjugation in homogeneous coordinates) is an orientation-preserving isometry that preserves the canonical Spin$^c$ structure. Combined with the identity on $S^3$, this defines a smooth CP action on $M$ preserving the Spin$^c$ structure and gauge bundle $E$.
Therefore $T_{\mathrm{CP}}$ is a closed Spin$^c$ 8-manifold, and by Theorem~\ref{thm:eta_even_vanishes}:
\begin{equation}
\eta(D_{T_{\mathrm{CP}}})=0 \in 4\mathbb{Z},
\qquad
A_{\mathrm{CP}}=\exp\!\left(\frac{i\pi}{2}\cdot 0\right)=1.
\end{equation}
\end{corollary}

\noindent\textbf{Remark.}
This result does \emph{not} depend on a delicate explicit evaluation of $\eta$;
it uses only the structural fact that the operator in Eq.~\eqref{eq:L_anomaly_phase} is a twisted Dirac operator on an even-dimensional closed manifold, hence has exact $\pm\lambda$ spectral pairing by Eq.~\eqref{eq:L_chirality_odd}.
For references stating this standard vanishing, see Loya--Moroianu--Park~\cite{LoyaMoroianuPark2009}.

\subsection{Main theorem: Strong CP solved}

\begin{theorem}[Strong CP all-loops closure]\label{thm:strong_cp_all_loops}
In the DFD microsector on $M = \mathbb{C}P^2\times S^3$ with the Standard Model fermion content:
\begin{enumerate}
\item The microscopic theory is CP invariant at $\theta_{\rm bare}=0$ (tree-level verified).
\item The CP anomaly phase is trivial: $A_{\mathrm{CP}}=1$ (Corollary~\ref{cor:DFD_ACP_trivial}).
\end{enumerate}
Therefore $\bar\theta=0$ to all loop orders. No axion is required.
\end{theorem}

\begin{proof}
Condition~(1) was established in Appendix~\ref{app:yukawa-hierarchy}: the K\"ahler structure ensures real Yukawa eigenvalues with $\arg\det(M_u M_d) < 10^{-19}$ rad.
Condition~(2) follows from Corollary~\ref{cor:DFD_ACP_trivial}: the mapping torus has even dimension (8), so the twisted Dirac operator has symmetric spectrum and $\eta = 0$ automatically.

Since both conditions hold, the renormalized effective action contains no CP-odd operators.
In particular, the coefficient of $\mathrm{Tr}(F\wedge F)$ vanishes identically at all scales.
\end{proof}

\subsection{Alternative verification: quaternionic structure}

An independent confirmation comes from the quaternionic structure on the $S^3$ factor.

\begin{lemma}[3D charge conjugation]\label{lem:S3_charge_conj}
Let $\sigma^a$ be Pauli matrices and consider the 3D Euclidean Dirac operator
$D_3 = i\sigma^a \nabla_a$.
Define the antiunitary charge conjugation $C_3\equiv \sigma^2\circ K$ (with $K$ complex conjugation).
Then
\begin{equation}
C_3^2=-1,\qquad C_3 D_3 C_3^{-1}=D_3 .
\end{equation}
\end{lemma}

\begin{proof}
The Pauli identity $\sigma^2(\sigma^a)^\ast\sigma^2=-\sigma^a$ implies
$C_3\,\sigma^a\,C_3^{-1}=-\sigma^a$, while antiunitarity gives $C_3\,i\,C_3^{-1}=-i$.
Therefore $C_3(i\sigma^a)C_3^{-1}=i\sigma^a$, proving commutation with $D_3$.
Finally $C_3^2=\sigma^2(\sigma^2)^\ast=-\mathbb{1}$.
\end{proof}

The quaternionic structure ($J^2=-1$) forces the fermion determinant to be real and nonnegative~\cite{FreedMoore2013,GarciaEtxebarriaMontero2018}, providing an independent confirmation that $A_{\mathrm{CP}}=1$.

\subsection{Falsifiable prediction}

Theorem~\ref{thm:strong_cp_all_loops} implies:
\begin{itemize}
\item \textbf{No QCD axion exists.} Axion searches (ADMX, ABRACADABRA, CASPEr, etc.) will find nothing.
\item \textbf{Any observed $\bar\theta \neq 0$ would falsify this mechanism.}
\end{itemize}
This is a sharp, experiment-confrontable prediction distinguishing DFD from Peccei--Quinn solutions.

\subsection{Summary: why the $S^3$ factor does quadruple duty}

The Strong CP problem is solved in DFD by topology, not by introducing new particles.
The key insight is dimensional: the microsector $M = \mathbb{C}P^2 \times S^3$ has $\dim M = 7$, so the mapping torus has $\dim T_{\mathrm{CP}} = 8$ (even), forcing $\eta = 0$ by spectral symmetry.

The same $S^3$ factor that:
\begin{enumerate}
\item Counts generations: $N_{\rm gen} = 3$ from the index theorem,
\item Stabilizes protons: baryon number is $\pi_3(S^3) = \mathbb{Z}$ winding,
\item Provides gauge emergence: $\pi_3(\mathrm{SU}(3)) = \mathbb{Z}$,
\end{enumerate}
also contributes the crucial ``$+1$'' to make $\dim T_{\mathrm{CP}} = 8$ even, thereby solving Strong CP.
This is a remarkable quadruple duty for one topological structure.
